\section{Conclusions}

The CCB test-beam analysis provides a controlled test of a segmented plastic-scintillator range stack for low-energy charged particles. Existing selected-data products indicate different longitudinal B-stack populations for the declared Sample-I and Sample-II run families, but the final quantitative depth-profile claim remains pending regeneration from the pre-threshold 8-channel-by-16-sample event population with run-aware controls. The analysed B readout samples only alternating longitudinal positions, so any range or energy-loss estimator is sensitive to where the stopping-power rise falls relative to instrumented positions. The intended telescope-style amplitude observable is \(\Delta E=A(B2)\) and \(E=A(B4)+A(B6)+A(B8)\); a B2--B4 plot is a two-channel diagnostic rather than the complete downstream-energy proxy.

The located 8-by-16 waveform product supports a negative timing conclusion rather than a detector-resolution number. In the current reconstructed subset, the B4--B6 pair residual has a central width of approximately 38 ns and is affected by the finite waveform representation, component-selection ambiguity and unresolved timing-reference/covariance terms. Historical sub-nanosecond simulation or incompatible-waveform results must therefore not be transferred to this beam-data product.

A single-stave Geant4 optical model remains valuable as a response framework, but the historical July calibration campaign is superseded as the nominal model and must be regenerated after source-level optical repairs. Likewise, the historical held-out energy reconstruction is non-authorising because its target was Birks-visible energy while the manuscript previously described raw deposited energy, its uncertainty/holdout coverage is inadequate for a resolution claim, and its underlying optical campaign is superseded.

The paper is therefore not yet a quantitative detector-performance publication. The shortest defensible path to submission is to close the source-bound hardware/run/polarity and 8-by-16 data contracts first, produce the final longitudinal beam-data result, establish a provenance-complete MC event measure, repair and rerun the \(\Delta E\)--\(E\) comparison, and then decide whether the optical/energy programme is mature enough to remain in the same paper. If it is retained, the optical grid must be regenerated with explicit \(E_{\rm raw}\) and \(E_{\rm vis}\) semantics, response-stage nuisances and a broader genuinely held-out reconstruction study. Every headline result should then be bound to its machine-readable source table, uncertainty, canonical claim status and exact reviewed source head before submission.
